\section*{September 7, 2026: Inspect three construction groups}
\textit{Agent-drafted evidence review; no new measurement decisions adopted.}

The evidence supports combining the two buildings at 4400 Grove, as Jacob
approved, but exposes a land mismatch. Two construction permits each identify
42 units, and the architect reports the first phase opening in October 2020.
The current density calculation divides these 84 units by 4.5 acres taken from
the architect's project page. The City's original approval instead assigns the
84-unit phase to a separate subarea with 102,637 square feet; the wider plan
contains other construction areas. The two assessed parcels have 97,215 square
feet, independently reproduced by their archived 2021 polygons. The smaller
difference between approved and assessed land still needs reconciliation. The
current whole-development denominator is unsupported for the two-building phase.

At Natchez, the City plan distinguishes a northern 72-unit area, a southern
84-unit area, and a western 39-unit expansion. Fourteen 2017 permits each
authorize six units in the southern area. Its archived parcel area agrees with
the 210,101-square-foot Assessor denominator. Keeping this defined group together
is supported more clearly than merging all construction with the same
development name. Completion remains uncertain: the 2018 plan labels the area
under construction, while the commercial Assessor records change its reported
construction year from 2017 to 2020 without changing floor or land area. The
available permit milestones establish subsequent completion, not its date.

The Roosevelt Square building at 1217 W Arthington has two additional conflicts.
Assessor reports from 2009 through 2020 give construction year 2005; reports from
2021 onward give 2006. All report six units, but an exact-address, exact-PIN
renovation permit from 2022 describes four. Neither discrepancy is resolved by
the larger development's Phase I completion date. Changing the building year to
2005 would put it outside the study period. Changing only its units to four would
reduce dwelling units per acre by one third. Neither change has been adopted.

\small
\textit{Evidence and scope.} Details, source links, record identifiers, and
arithmetic scenarios are in the release audit's
\path{phased_project_review.md}. Its Make target produces 181 permit search
records with unique source IDs in \path{phased_project_permits.csv}; these are
evidence candidates, not accepted project matches. The new public PD 1395 download
has a concrete acquisition rule and recorded checksum. The logbook build holds
three unrelated existing audit products fixed. Production datasets and paper
estimates are unchanged; this is not an end-to-end replication result.
\normalsize
